% !TeX root = ../main.tex
\section{Discussion}
\label{sec:discussion}

\subsection{Separating Algorithmic and Physical Frontiers}
\label{subsec:separating-frontiers}

The primary finding of this behavioral simulation study is the clear separation between algorithmic generalization risks and physical device limits in analog CIM deployment. Standard fixed-mask hardware-aware training, while effective at higher precisions, fails fundamentally at extreme low precision (4-bit) because the network overfits to a single hardware instance's mismatch field. Ensemble HAT resolves this by treating hardware variability as a distribution-level regularizer during training, effectively "rescuing" the 4-bit regime for hardware instances unseen during training.

This rescue, however, does not eliminate physical constraints; it merely reveals them. Once algorithmic convergence is secured, the system is bounded by the precision-retention frontier of the underlying memory technology. In our tested PCM UnitCell models, 8-bit precision emerges as the strongest practical operating point, combining high fresh accuracy (77.60\%) with drift stability. 6-bit precision, while drift-stable, enters a D2D-sensitive regime where device mismatch degrades mean accuracy (68.55\%) and introduces high seed-to-seed variance (std 6.03\%)---a transition zone between the quantization-dominated 4-bit regime and the noise-absorbed 8-bit regime.

\subsection{Hard Cliffs and Operational Envelopes}
\label{subsec:hard-cliffs}

The joint ADC--D2D sensitivity analysis identifies a hierarchical set of hardware requirements. ADC resolution serves as a "hard gate": below 6 bits, classification accuracy collapses regardless of device quality. Above this threshold, device-to-device variability becomes the primary determinant of accuracy. This hierarchy suggests a clear hardware roadmap for vision transformer deployment: prioritize a high-precision (6-bit or greater) readout interface before attempting to minimize device-level mismatch.

\subsection{Limitations and outlook}
\label{subsec:limitations}

\paragraph{Scope.} The conclusions are bounded by the present behavioral simulation stack. Device profiles are calibrated from literature values or a single fabrication batch, so the numbers do not claim cross-batch, temperature, aging, or circuit-layout coverage. We model quantization, D2D/C2C variability, ADC precision, write nonlinearity, and retention, but defer parasitic IR drop, sneak paths, and pulse-faithful programming to follow-on hardware-calibrated studies.

\paragraph{Energy and scale.} The energy estimates are first-order analytical projections from literature-proxy constants ($E_{\text{cell}}=100$~fJ, $E_{\text{ADC},8b}=25$~fJ, $E_{\text{DAC},8b}=30$~fJ; Section~\ref{subsec:sensitivity-energy}). They rank deployment risks rather than predict silicon performance. The current ablation matrix also prioritizes non-ideality coverage over dataset breadth; cross-architecture TinyImageNet sweeps (see Supplementary Information) are treated as validation, not as main-text claims.

\begin{table}[htbp]
\centering
\caption{\textbf{Deployment Risk Ranking and Mitigation Strategies.} Factors are ranked by their impact on categorical accuracy as identified by Sobol sensitivity analysis and precision-ladder sweeps.}
\label{tab:deployment_risks}
\begin{tabular}{lll}
\toprule
\textbf{Factor} & \textbf{Impact Rank} & \textbf{Mitigation Strategy} \\
\midrule
ADC Resolution & High (Hard Gate) & Secure $\geq$6-bit readout precision \\
Stochastic D2D & High (Generalization) & Ensemble HAT (distribution matching) \\
PCM Drift & Moderate (Time-bound) & 8-bit for stability; refresh 4-bit periodically \\
Write Nonlinearity & Moderate (Optimization) & STE gradient scaling \\
Shot Noise & Low (Transduction) & Front-end inverse-gamma compensation \\
\bottomrule
\end{tabular}
\end{table}

\paragraph{Outlook.} The next tests are straightforward: replace UnitCell presets with measured-array statistics, verify whether the non-monotonic precision-accuracy curve and 6-bit transition-zone behavior persist under measured-device statistics, and evaluate memory-bound inference components such as analog KV-cache with explicit train/eval noise separation. These extensions do not alter the present conclusion: algorithmic cross-instance resampling and physical precision-retention selection are separate deployment decisions.
